\section{Estratégia de aproximação relativa}
\label{sec:estrategia_orbital_global}

Ao final da manobra de transferência orbital, o perseguidor é inserido em uma nova órbita e permanece em regime de \textit{drift} por um intervalo de tempo predeterminado.
Nesse intervalo, a dinâmica é propagada essencialmente pelo problema de dois corpos, enquanto são realizadas checagens de posição e velocidade relativa em relação ao alvo.

\subsection{Aproximação curta}
\label{sec:aproximacao_curta}

Somente após essa fase de \textit{drift} é que se considera o início da aproximação curta.
O estado do perseguidor ao término do \textit{drift} é tomado como condição inicial no referencial \gls{lvlh} do alvo, e a dinâmica relativa passa a ser descrita pelo modelo linearizado \gls{hcw}.
A partir desse instante, a lei de controle \gls{pd} assume realizando correções nas componentes radial, ao longo da órbita e transversal.

\subsection{Aproximação final}
\label{sec:aproximacao_final}

A aproximação final corresponde ao trecho final da manobra translacional, já no interior da região de captura em \gls{lvlh} centrado no alvo.
Nessa fase, o perseguidor encontra-se a poucos metros da porta de acoplamento e as condições de posição, velocidade relativa e atitude são mais restritivas do que nas etapas anteriores.
O objetivo principal é reduzir os erros residuais de posição e velocidade até atingir um estado quase estacionário em relação ao alvo para que ocorra a atuação do manipulador robótico.

Assim que o perseguidor insere o efetuador final na porta de acoplamento, a aproximação final é finalizada, assim como toda a missão.
A Tabela~\ref{tab:parametros_pos_TH} demonstra os valores obtidos durante a aproximação curta e final.

\begin{table}[H]
\centering
\renewcommand{\arraystretch}{1.2}
\caption{Parâmetros relativos alvo--perseguidor durante a aproximação curta.}
\label{tab:parametros_pos_TH}
\begin{tabular}{lrr}
\hline
\textbf{Variável} & \textbf{Valor} & \textbf{Unidade} \\
\hline
\multicolumn{3}{l}{\textbf{Condições relativas após TH}} \\
\hline
$\Delta r_{\text{radial,TH}}$& $-100.00$& m \\
$\Delta r_{\text{along,TH}}$& $-500.00$& m \\
$\Delta r_{\text{cross,TH}}$& $0.00$& m \\
\hline
\multicolumn{3}{l}{\textbf{Limites de velocidade relativa}} \\
\hline
$v_{\text{rel,max}}(d \leq 1000\,\text{m})$ & $0.30$ & m/s \\
$v_{\text{rel,max}}(d \leq 50\,\text{m})$& $0.03$ & m/s \\
\hline
\multicolumn{3}{l}{\textbf{Condições de atitude inicial}} \\
\hline
$q_{\text{alvo},0}$& $[1,\;0,\;0,\;0]$ & -- \\
$\boldsymbol{\phi}_{\text{chr},0}$ (Euler) & $[7,\;8,\;9]$     & deg \\
$q_{\text{chr,des}}$& $[1,\;0,\;0,\;0]$ & -- \\
\hline
\end{tabular}
\end{table}

A Tabela~\ref{tab:aproximacao_final} explicita os valores obtidos durante a aproximação final.
\begin{table}[H]
\centering
\renewcommand{\arraystretch}{1.2}
\caption{Condições de referência na aproximação final em LVLH.}
\label{tab:aproximacao_final}
\begin{tabular}{lrr}
\hline
\textbf{Variável} & \textbf{Valor} & \textbf{Unidade} \\
\hline
\multicolumn{3}{l}{\textbf{Posição em LVLH (alvo como referência)}} \\
\hline
Posição do alvo em LVLH, $\vec r_{\text{alvo}}$ &
$[0.00,\;0.00,\;0.00]$ & m \\
Posição do COM da base em LVLH, $\vec r_{\text{base}}$ &
$[0.00,\;0.00,\;0.70]$ & m \\
\hline
\multicolumn{3}{l}{\textbf{Velocidade relativa na aproximação final}} \\
\hline
Velocidade relativa base–alvo, $\dot{\vec r}_{\text{rel}}$ &
$[0.00,\;0.00,\;0.00]$ & m/s \\
\hline
\end{tabular}
\end{table}
